\chapter{Formulação do Princípio}
\label{chap:formulacao}

Seja $R$ o contrato semântico confirmado, $P$ o processo de referência e
$T_R(P)$ o processo governado por uma política \LangT. A equivalência
operacional é relativa aos observáveis de $R$:

\begin{equation}
  T_R(P) \equiv_R P
  \quad\Longleftrightarrow\quad
  \forall r \in R_{\mathrm{HARD}},\;
  V_r\!\left(T_R(P)\right)=V_r(P)=\mathrm{PASS}.
  \label{eq:equivalencia}
\end{equation}

Essa expressão não implica identidade de bits ou de estado intermediário. A
compressão em uma dimensão $k$ existe somente se a contabilidade congelada
mostrar economia líquida:

\begin{equation}
  C_k\!\left(T_R(P)\right)
  = C_k(P)-C_k(E_{\mathrm{redundante}})+C_k(T),
  \qquad C_k(E_{\mathrm{redundante}})>C_k(T).
  \label{eq:compressao}
\end{equation}

\section{Cadeia de autoridade}

\begin{figure}[htbp]
\centering
\begin{tikzpicture}[
  node distance=7mm and 5mm,
  box/.style={draw=CasuloNavy,rounded corners,fill=CasuloMist,
              align=center,minimum height=9mm,text width=23mm,font=\small\sffamily},
  arrow/.style={-{Stealth[length=2mm]},thick,color=CasuloTeal}
]
\node[box] (intent) {Intenção};
\node[box,right=of intent] (contract) {Contrato $R$};
\node[box,right=of contract] (t) {Política \LangT};
\node[box,right=of t] (tcap) {IR / \TCAP};
\node[box,below=of tcap] (primitive) {Primitive};
\node[box,left=of primitive] (receipt) {Receipts};
\node[box,left=of receipt] (judge) {Juiz \LangT};
\node[box,left=of judge] (publish) {Publicação};
\draw[arrow] (intent)--(contract);
\draw[arrow] (contract)--(t);
\draw[arrow] (t)--(tcap);
\draw[arrow] (tcap)--(primitive);
\draw[arrow] (primitive)--(receipt);
\draw[arrow] (receipt)--(judge);
\draw[arrow] (judge)--(publish);
\end{tikzpicture}
\caption{Cadeia conceitual de autoridade. O manuscrito final deve ligar cada
nó a símbolos e artefatos reais.}
\label{fig:cadeia}
\end{figure}
